\documentclass[12pt]{iopjournal}

\usepackage{bm}
\usepackage{amssymb}
\usepackage{booktabs}
\usepackage{mathtools}
\usepackage{microtype}
\usepackage{url}
\usepackage{placeins}
\usepackage[backend=biber,style=numeric-comp,sorting=none,giveninits=true,maxbibnames=99]{biblatex}
\addbibresource{references.bib}

\hypersetup{
  colorlinks=true,
  allcolors=blue,
  pdftitle={Supplementary Material for Detector-Restricted Identifiability of Charged Spin-1/2 and Spin-1 Particles},
  pdfauthor={Dino Vlahek}
}

\newcommand{\DKL}{D_{\mathrm{KL}}}
\newcommand{\Tr}{\operatorname{Tr}}
\newcommand{\dd}{\mathrm d}
\providecommand{\alttext}[1]{}
\newcommand{\sfig}[2]{\includegraphics[width=0.88\linewidth]{figures/#1}\alttext{#2}}
\newenvironment{ruledtabular}{\centering}{}
\newcommand{\colrule}{\hline}

\renewcommand{\theequation}{S\arabic{equation}}
\renewcommand{\thefigure}{S\arabic{figure}}
\renewcommand{\thetable}{S\Roman{table}}

\setlength{\headheight}{15pt}
% Author-prepared review PDF: replace template-only production headers.
\renewcommand{\articletype}[1]{\par\noindent\textsf{\textbf{#1}}\par\medskip}
\AtBeginDocument{%
  \fancyhfoffset[L]{0pt}
  \fancyhead[L]{\small Detector-restricted spin identifiability}
  \fancyhead[R]{\small D. Vlahek}
  \renewcommand{\headrulewidth}{0.3pt}
}

\begin{document}
\articletype{Supplementary Material}
\title{Supplementary Results}
\author{Dino Vlahek}\par
\affil{Laboratory for Data Mining and Intelligent Systems, Faculty of Organization and Informatics, University of Zagreb, Pavlinska 2, 42000 Vara\v{z}din, Croatia}
\par\email{dvlahek@foi.hr}\par

\begin{abstract}
Numerical protocols, nuisance-class robustness results, whole-class certificates, finite-basis audits, and estimator-space checks supporting the main article.\par
\medskip
\noindent\textbf{Associated article:} ``Detector-Restricted Identifiability of Charged Spin-$1/2$ and Spin-$1$ Particles: Whole-Class Certificates and Exact Degeneracies.''\par
\medskip
All claims about detector-restricted identification refer to the model class and measurement architecture specified in the main article.
\end{abstract}

% ============================================================
\section{Finite-bandwidth validation protocol}
% ============================================================

The finite-width calculation uses the dimensionless transfer variable
\begin{equation}
x=Q^2/m^2
\end{equation}
and the $\delta=0$ reference-slice conditions
\begin{equation}
G_1(0)=1,\qquad G_M(0)=g_M,\qquad G_Q(0)=1-g_M.
\end{equation}
The robustness audit uses $G_Q(0)=1+\delta-g_M$ with $|\delta|\le\Delta$.
The admissible class is defined for every $x,y\ge0$ by
\begin{equation}
|G_i(x)-G_i(y)|\le L|x-y|,
\qquad i\in\{1,M,Q\}.
\end{equation}
The equivalent bound $|G_i'|\le L$ holds almost everywhere for the piecewise-linear constructions. The principal bounds are defined directly on the full anchored Lipschitz class.

For $r=k/m$, $\widehat\sigma=\sigma_k/m$, $\omega_r=\sqrt{1+r^2}$, and $x=2(\omega_r-1)$, the common Gaussian window defined in the main article and the squared current kinematics induce
\begin{align}
\rho_0(r;\widehat\sigma)&=
\frac{e^{-r^2/(2\widehat\sigma^2)}}{2\omega_r},\\
\dd\nu_M&=\frac{\rho_0 r^2\,\dd r}{\int\rho_0 r^2\,\dd r},\\
\dd\nu_C&=\frac{\rho_0r^2x\,\dd r}
{\int\rho_0r^2x\,\dd r}.
\end{align}
The charge weight $\rho_0r^2x$ follows from the squared longitudinal-current rank-$2$ response. Both channels use the common Gaussian window. If $N_M$ and $N_C$ denote the two denominators above, the channel calibration, after the fixed rank-$2$ constants are removed, is $\widetilde W_a=N_a\overline W_a$. The charge coordinate is formed only after division of the Breit-frame rank-$2$ matrix coefficient by its known factor $2\eta/3=x/6$. These statements fix the relative normalization used before any covariance whitening. The displayed measures are the exact measures used in every finite-bandwidth table. Their numerical normalizations equal one to better than $3\times10^{-16}$ on the production grid. At $\widehat\sigma=0.15$, $\langle x\rangle_C=0.1054292265$.

We examine the effect of an additional momentum-dependent charge acceptance
\begin{equation}
A_C^{(\alpha)}(k)=\left[\frac{2m}{E_k+m}\right]^\alpha,
\qquad \alpha\in\{0,1,2,3\},
\label{eq:acceptance-family}
\end{equation}
The main model uses $\alpha=0$. Positive $\alpha$ values define alternative charge-sector acceptances. For each $\alpha$ we recompute the pointwise full-Lipschitz charge envelope and evaluate the same two-kink witness family at every grid cell. Table~\ref{tab:acceptance-sensitivity} reports the resulting \emph{sufficient} certificates and constructive overlaps. No physical prior on $\alpha$ is inferred from this illustrative sensitivity calculation.

The same 33 two-kink witnesses retain their endpoint sign changes across all four tested choices. For $\alpha=0$ and $1$, the boundary cell $(\sigma_k/m,L)=(0.50,1)$ remains unresolved. It becomes certified for $\alpha=2$ and $3$. The magnetic-only lower bound and the static primitive-resource criterion are unchanged. Other selections or uncertainty sets for $\alpha$ require separate analysis. The four tested choices do not establish a uniform charge-channel result.

\begin{table}[h]
\caption{Acceptance-sensitivity audit for the 99-point $(\sigma_k/m,L)$ grid. A certified point has a positive full-Lipschitz envelope. An overlap has a signed two-kink witness. A nonpositive envelope without a witness is unresolved, not proof of degeneracy. The counts are for each fixed $\alpha$ separately, not one certificate uniform in $\alpha$.}
\label{tab:acceptance-sensitivity}
\begin{ruledtabular}
\begin{tabular}{r r r r r}
$\alpha$ & certified & overlap & unresolved & $B_C(0.50,1)$ \\
\colrule
0 & 63 & 33 & 3 & $-0.111794$ \\
1 & 63 & 33 & 3 & $-0.051924$ \\
2 & 64 & 33 & 2 & $+0.002340$ \\
3 & 64 & 33 & 2 & $+0.051544$ \\
\end{tabular}
\end{ruledtabular}
\end{table}

Componentwise moment bounds certify 46 points. A positivity reduction adds 17 whole-class certificates. The two-kink family yields 33 overlaps. Finite-basis QCQP and SOS calculations provide numerical checks, and Sec.~\ref{sec:supp-K-sensitivity} gives the dependence on the magnetic-moment domain.

The high-resolution validation configuration is summarized in Table~\ref{tab:scan-config}.

\begin{table}[h]
\caption{Finite-bandwidth high-resolution validation configuration.}
\label{tab:scan-config}
\begin{ruledtabular}
\begin{tabular}{ll}
quantity & value \\
\colrule
$g_M$ interval & $[-4,4]$ \\
$\sigma_k/m$ & $0.02,0.05,0.08,0.10,0.15,0.20,0.30,0.40,0.50$ \\
$L$ & $0.25,0.5,1,2,3,4,6,8,12,16,24$ \\
scan points & $9\times11=99$ \\
global QCQP basis & 8 segments \\
adversarial multistart basis & 48 segments, 200 random starts \\
finite-basis spectral grid & 2601 points \\
finite-basis spectral cutoff & $9\sigma_k$ \\
whole-class envelope grid & 20001 points to $12\sigma_k$ \\
global QCQP audit & $2^8=256$ convex QPs per point \\
explicit-overlap construction & 32 segments \\
SOS relaxation & verified degree 2 \\
Monte Carlo draws per hypothesis & 8000 \\
model-cloud size & 30000 \\
$g$ sampling in inference cloud & independent $\mathrm{Unif}[-4,4]$ \\
form-factor slope sampling & 24 independent $\mathrm{Unif}[-L,L]$ draws \\
independent truth-cloud size & 10000 \\
\end{tabular}
\end{ruledtabular}
\end{table}

For the unit-normalized dipole proxy $\Lambda/m=\sqrt{2/L}$, the values
corresponding to $L=(0.25,0.5,1,2,3,4)$ are
$(2.828,2.000,1.414,1.000,0.816,0.707)$, respectively.  For
$L=(6,8,12,16,24)$ they are
$(0.577,0.500,0.408,0.354,0.289)$.

The proxy row follows $\Lambda/m=\sqrt{2/L}$ for a unit-normalized dipole that saturates the origin slope bound. It supplies local scale intuition. For a form factor with a nonunit anchor, the absolute origin slope also contains the anchor magnitude. The Lipschitz class leaves the detailed dipole form and structural scale free.

\subsection{Componentwise moment bound}

Let $\langle x^n\rangle_a=\int\dd\nu_a(x)x^n$. A direct componentwise error box obeys
\begin{equation}
|\overline W_M-g^2|
\le2|g|L_M\langle x\rangle_M+L_M^2\langle x^2\rangle_M,
\end{equation}
and
\begin{align}
|\overline W_C-(1-g)|\le{}&
[L_Q+|1-g|L_1]\langle x\rangle_C
+L_1L_Q\langle x^2\rangle_C
\nonumber\\
&+\frac14(1-g)^2\langle x\rangle_C
+\frac12|1-g|L_Q\langle x^2\rangle_C
+\frac14L_Q^2\langle x^3\rangle_C.
\end{align}
These inequalities certify 46 grid points. The positivity envelope below certifies 17 additional points and gives 63 whole-class certificates in total.

\subsection{Whole-Lipschitz positivity reduction}

The magnetic response is an integral of a nonnegative function against a positive measure whose closed support contains the anchor $x=0$,
\begin{equation}
\overline W_M=\int d\nu_M\,G_M^2.
\end{equation}
Therefore continuity and $0\in\overline{\operatorname{supp}\nu_M}$ imply that $\overline W_M=0$ requires $g_M=0$ and $G_M\equiv0$ on the detector support. The static anchors in the remaining charge sector are then $G_1(0)=G_Q(0)=1$. At each transfer $x$, both form factors lie in $[1-Lx,1+Lx]$. The pointwise lower envelope is
\begin{equation}
h_L(x)=\min_{a,b\in[1-Lx,1+Lx]}
\left(ab+\frac{x}{4}b^2\right),
\end{equation}
and the whole-class lower bound is
\begin{equation}
B_C(\sigma_k/m,L)=\int d\nu_C(x)\,h_L(x).
\end{equation}
The inner minimization is deterministic. At fixed $b$ the minimum over $a$ is on an interval endpoint. For either endpoint the objective is a convex quadratic in $b$, whose stationary point $-2a/x$ is clipped to the allowed interval. Under $\overline W_M=0$, every admissible profile obeys $\overline W_C\ge B_C$. Thus $B_C>0$ excludes a joint origin overlap in the complete anchored Lipschitz class, without a piecewise-linear truncation.

The pointwise minimization relaxes the anchored functional problem because the minimizing values at different $x$ need not form one common continuous Lipschitz pair anchored at unity. The resulting analytic lower bound applies to the complete anchored Lipschitz class and can be conservative. A value $B_C\le0$ leaves the point undecided. This occurs at the three points retained as unresolved even though their finite piecewise-linear audit minima are positive.

This bound certifies 17 of the 21 audited cells left inconclusive by the componentwise error box. The final 99-point count is 63 certified, 33 explicit overlaps, and 3 unresolved. The remaining points are $(\sigma_k/m,L)=(0.08,24)$, $(0.10,16)$, and $(0.50,1)$.

\subsection{Whole-class distance enclosure}

For a general anchor $g$, the anchored Lipschitz inequalities imply
\begin{align}
\overline W_M&\ge M_L(g)
=\int\dd\nu_M[\max(|g|-Lx,0)]^2,\\
C_L^-(g)&\le\overline W_C\le C_L^+(g),
\end{align}
where $C_L^\pm$ are the integrals of the pointwise extrema of
\begin{equation}
ab+\frac{x}{4}b^2,\qquad
a\in[1-Lx,1+Lx],\quad
b\in[1-g-Lx,1-g+Lx].
\end{equation}
Thus the whole function class at fixed $g$ is contained in the half-strip
$[M_L(g),\infty)\times[C_L^-(g),C_L^+(g)]$. Minimizing the covariance-whitened norm over this half-strip and then over $g\in[-4,4]$ gives a valid lower bound on the distance of the complete class. This construction gives the bound reported in Table~\ref{tab:covariance}.

At the reference point $(\sigma_k/m,L)=(0.10,2)$ and unit covariance shape, the analytic whole-class lower bound is $0.4034833471$ at $g=0.5735325$. The explicit admissible profiles $G_M=\operatorname{sgn}(g)\max(|g|-Lx,0)$, $G_1=1-Lx$, and $G_Q=1-g-Lx$ give the upper bound $0.4034946602$ at $g=0.5734946$. Thus
\begin{equation}
0.4034833471\le d_{\rm class}\le0.4034946602.
\end{equation}
The bracket width is $1.13131\times10^{-5}$. The reference integral uses a symmetric grid of $20001$ points truncated at $12\sigma_k$. The half-strip inequality is analytic, while its numerical evaluation is checked by grid and Gaussian-tail convergence. The printed decimal values are not interval-arithmetic bounds.

The same deterministic minimization was evaluated at all 99 bandwidth--smoothness points. Figure~\ref{fig:supp-distance-grid} shows the whole-class unit-covariance lower bound. Exactly 63 values are positive above $10^{-10}$, matching the binary whole-class certificate count. The 33 explicit overlaps and the 3 unresolved points have a zero enclosure bound.

\begin{figure}[t]
\centering
\sfig{figureS1.png}{Quantitative whole-class distance lower bound on the complete bandwidth--smoothness grid.}
\caption{Whole-class unit-covariance lower bound $D_L(I)$ on all 99 points. The color scale is logarithmic. Gray indicates that the half-strip enclosure reaches the origin and supplies no positive distance certificate. Explicit overlaps and unresolved points are distinguished by the separate constructive search.}
\label{fig:supp-distance-grid}
\end{figure}

\subsection{Quadrature and Gaussian-tail convergence}

Table~\ref{tab:alpha0-convergence} reports the reference lower-bound calculation on three grid resolutions and three spectral cutoffs. The difference between the largest and smallest values is approximately $6\times10^{-11}$. The minimizer $g$ remains close to $0.57353$. These numerical checks support the reported decimal precision but do not constitute interval-arithmetic certification.

\begin{table}[ht]
\caption{Convergence of the whole-class distance lower bound at $(\sigma_k/m,L)=(0.10,2)$ for unit covariance.}
\label{tab:alpha0-convergence}
\begin{ruledtabular}
\begin{tabular}{ccc}
cutoff & grid points & $D_L(I)$ \\
\colrule
$10\sigma_k$ & 10001 & 0.403483347114 \\
$10\sigma_k$ & 20001 & 0.403483347140 \\
$10\sigma_k$ & 40001 & 0.403483347138 \\
$12\sigma_k$ & 10001 & 0.403483347112 \\
$12\sigma_k$ & 20001 & 0.403483347132 \\
$12\sigma_k$ & 40001 & 0.403483347138 \\
$14\sigma_k$ & 10001 & 0.403483347086 \\
$14\sigma_k$ & 20001 & 0.403483347146 \\
$14\sigma_k$ & 40001 & 0.403483347138 \\
\end{tabular}
\end{ruledtabular}
\end{table}

\subsection{Coordinate convention audit}

The primary nuisance class bounds the covariant vertex variables $(G_1,G_M,G_Q)$. Since $G_1=G_C-xG_Q/6$, the charge response is equivalently
\begin{equation}
\overline W_C=\int\dd\nu_C\,
G_Q\left(G_C+\frac{x}{12}G_Q\right).
\end{equation}
Imposing the same numerical $L$ on $(G_C,G_M,G_Q)$ defines a different function class. Repeating the exact-origin envelope replaces $ab+xb^2/4$ by $ab+xb^2/12$. Because $x\ge0$, the latter integrand is no larger for the same $(a,b)$ box, so the Sachs test cannot certify more points by construction. Figure~\ref{fig:supp-coordinate-audit} compares the two sufficient tests. They exclude the origin at 63 and 62 of the 99 grid points, respectively. Only $(0.30,2)$ changes. Its lower envelopes are $+0.008404$ in the covariant class and $-0.005405$ in the Sachs class. At the reference point the corresponding whole-class distance lower bounds are $0.403483$ and $0.402935$.

\begin{figure}[h]
\centering
\sfig{figureS2.png}{Comparison of sufficient whole-class origin-exclusion tests in covariant and Sachs Lipschitz coordinates.}
\caption{Coordinate-convention audit. Green denotes a positive pointwise charge envelope after magnetic positivity reduction. Gray marks points left undecided by this sufficient criterion. The two panels use different smoothness classes.}
\label{fig:supp-coordinate-audit}
\end{figure}

\subsection{Dipole scale and local finite-transfer stability}

A unit-normalized dipole profile provides a local interpretation of the Lipschitz bound,
\begin{equation}
\widehat G(Q^2)=\left(1+\frac{Q^2}{\Lambda^2}\right)^{-2},
\qquad
\left|\frac{\mathrm d\widehat G}{\mathrm dx}(0)\right|=\frac{2m^2}{\Lambda^2},
\qquad \frac{\Lambda}{m}=\sqrt{\frac{2}{L}}.
\end{equation}
The last equality holds when the origin slope saturates $L$. It gives $\Lambda=m$ at $L=2$ and $\Lambda\simeq0.289m$ at $L=24$. For an anchored profile $G_i(Q^2)=G_i(0)\widehat G(Q^2)$, the absolute slope also depends on $|G_i(0)|$. This dipole example is illustrative and is not imposed in the whole-class analysis.

There is also a simple stability statement for the local, amplitude-calibrated protected coordinates. Let
\begin{equation}
\bm W_1(g,Q^2)=\bigl(|G_M(g,Q^2)|^2,G_Q(g,Q^2)\bigr)
\end{equation}
be jointly continuous on $K\times[0,Q_{\max}^2]$ for compact $K$, and let $\|\bm v\|_A=\sqrt{\bm v^{\mathsf T}A\bm v}$ with positive-definite $A$. Suppose
\begin{equation}
d_0=\min_{g\in K}\|\bm W_1(g,0)\|_A>0
\end{equation}
and that for some $Q_\ast^2>0$,
\begin{equation}
\sup_{g\in K}\|\bm W_1(g,Q^2)-\bm W_1(g,0)\|_A<d_0
\qquad (0\le Q^2\le Q_\ast^2).
\end{equation}
The reverse triangle inequality gives
\begin{equation}
\|\bm W_1(g,Q^2)\|_A\ge d_0-\|\bm W_1(g,Q^2)-\bm W_1(g,0)\|_A>0.
\end{equation}
Hence the local response remains separated from the spin-$1/2$ protected origin throughout that interval. This continuity argument concerns local amplitude-calibrated coordinates, not the bilinear finite-width charge effect. The main-text whole-Lipschitz bounds establish the latter independently.


\subsection{Explicit continuous-class overlap}

A single deterministic two-kink family produces a witness for each of the 33 overlap cells. For each width $\widehat\sigma=\sigma_k/m$ set
\begin{align}
X&=2\left(\sqrt{1+(9\widehat\sigma)^2}-1\right),
&h&=X/32,\notag\\
G_M^{(t)}(x)&=0,
&G_Q^{(t)}(x)&=1-tL\min(x,X),\notag\\
G_1^{(t)}(x)&=1+tL[\min(x,X)-2\min(x,h)],
&0&\le t\le1.
\end{align}
The $G_1$ slopes are $-tL$ on $(0,h)$, $+tL$ on $(h,X)$ and zero beyond $X$. The $G_Q$ slope is $-tL$ on $(0,X)$ and zero beyond. Thus each profile belongs to the complete anchored Lipschitz class for every $t\in[0,1]$. For the 33 overlap cells, direct continuous quadrature split at $h$ and $X$ gives $\overline W_C(0)>0$ and $\overline W_C(1)<0$. The smallest absolute negative endpoint magnitude on this set is $0.0474608$, attained at $(\widehat\sigma,L)=(0.15,8)$. Consequently, the intermediate-value theorem supplies an exact root for each of these 33 cells. The accompanying code and CSV include the endpoint pairs and numerical roots. Sign checks use numerical quadrature without formal interval certification.

At $(\sigma_k/m,L)=(0.15,8)$, the same family is represented by 32 equal-width segments for $G_1$ and $G_Q$. Scaling its slopes by $t\in[0,1]$ gives
\begin{equation}
\overline W_C(0)=1.0263573066,
\qquad
\overline W_C(1)\simeq-0.0474607720.
\end{equation}
The intermediate-value theorem therefore gives an exact root. Numerically,
\begin{equation}
t_\star\simeq0.9667494646,
\qquad
|\overline W_C(t_\star)|\sim10^{-15}\quad\text{(unrounded numerical root)},
\end{equation}
using knot-split continuous quadrature. The maximum absolute covariant slope at the root is $7.7339957<8$. Its exact existence follows from the endpoint sign change, not from floating-point rounding. Transforming the same witness with $G_C=G_1+xG_Q/6$ gives
\begin{equation}
\max|G_C'|\simeq7.791093<8,
\qquad
\max|G_Q'|\simeq7.733996<8.
\end{equation}
The displayed overlap therefore belongs to both the covariant and Sachs Lipschitz classes. The other 32 witnesses refer to the covariant class. The machine-readable audit contains the endpoint values and numerical roots for all 33 cases.

\subsection{Higher-resolution search and unresolved cells}

The pointwise envelope and the two-kink construction leave three cells unresolved. Their global eight-segment QCQP minima are listed in Table~\ref{tab:unresolved-searches}. We additionally optimize $48$-segment profiles with alternating minimization and $200$ seeded starts for each cell. The same procedure is applied to the known-overlap case $(0.15,8)$ as a positive control. Table~\ref{tab:highres-search} reports the best values after independent spectral-grid evaluation. A negative value on the control verifies that the search can recover a signed overlap in this setting. The three unresolved cases remain positive in this search. These are local multistart results and do not establish whole-class separation.

\begin{table}[ht]
\caption{Global eight-segment QCQP minima for the three unresolved cells. Positive minima apply only to this finite basis.}
\label{tab:unresolved-searches}
\begin{ruledtabular}
\begin{tabular}{ccc}
$\sigma_k/m$ & $L$ & $q_{8,\min}$ \\
\colrule
0.08 & 24 & 0.201665 \\
0.10 & 16 & 0.168785 \\
0.50 & 1 & 0.105516 \\
\end{tabular}
\end{ruledtabular}
\end{table}

\begin{table}[ht]
\caption{Adversarial search with 48 equal-width segments and 200 random starts per cell. The values are evaluated on an independent 20001-point spectral grid. The first row is a known-overlap control. Positive entries do not certify separation.}
\label{tab:highres-search}
\begin{ruledtabular}
\begin{tabular}{cccc}
$\sigma_k/m$ & $L$ & $\overline W_{C,48}$ & result \\
\colrule
0.15 & 8 & $-0.050566$ & overlap control \\
0.08 & 24 & $+0.064931$ & unresolved \\
0.10 & 16 & $+0.036632$ & unresolved \\
0.50 & 1 & $+0.044869$ & unresolved \\
\end{tabular}
\end{ruledtabular}
\end{table}

The slopes satisfy $|u_i|,|v_i|\le L$. The search starts from zero and extremal profiles and uses independent uniform random initial slopes. The spectral cutoff is $9\sigma_k$ for optimization and $12\sigma_k$ for re-evaluation. Profile arrays, seeds and evaluation results are supplied with the audit code.

\subsection{Global finite-basis QCQP and degree-$2$ SOS audit}

For integrated-slope basis functions $\bm\phi(x)$, write
\begin{equation}
G_1(x)=1+\bm\phi(x)^{\mathsf T}\bm u,
\qquad
G_Q(x)=1+\bm\phi(x)^{\mathsf T}\bm v.
\end{equation}
The charge response is the QCQP
\begin{equation}
\overline W_C=c_0+\bm a_u^{\mathsf T}\bm u
+\bm a_v^{\mathsf T}\bm v
+\bm u^{\mathsf T}H\bm v
+\bm v^{\mathsf T}Q\bm v,
\qquad Q\succeq0,
\end{equation}
with $|u_i|,|v_i|\le L$. For fixed $\bm v$ the objective is linear in $\bm u$, so a global minimum has $\bm u$ at one of the $2^8$ box vertices. At each vertex the remaining problem in $\bm v$ is a convex box QP. All 21 points inconclusive under the componentwise test have positive global minima in this eight-segment basis. The box-QP stationary conditions are checked numerically for each retained minimizer.

The separate degree-$2$ relaxation uses $\bm z=(\bm u,\bm v)$ and verifies
\begin{equation}
q(\bm z)-\gamma=
\begin{bmatrix}1\\\bm z\end{bmatrix}^{\!\mathsf T}
S
\begin{bmatrix}1\\\bm z\end{bmatrix}
+\sum_i\lambda_i(L^2-z_i^2),
\end{equation}
with $S\succeq0$ and $\lambda_i\ge0$. The degree-$2$ SOS calculation gives 12 positive numerical lower bounds among the 21 QCQPs. The remaining relaxation values are nonpositive and inconclusive. The accompanying data include multipliers and PSD eigenvalues for direct checking. The calculation uses floating-point arithmetic. The classification is based on whole-class bounds and constructive overlaps.

Table~\ref{tab:qcqp-sos} lists the complete 21-point audit and makes the difference between the pointwise whole-class relaxation and the finite-basis global minimum explicit.

\begin{table}[!htbp]
\caption{Independent audits of the 21 points inconclusive under the componentwise error box. $B_C$ is the whole-Lipschitz lower bound, $q_{8,\min}$ is the global eight-segment QCQP minimum, and $\gamma_2$ is the verified degree-$2$ SOS bound. A negative $\gamma_2$ is inconclusive.}
\label{tab:qcqp-sos}
\begin{ruledtabular}
\begin{tabular}{rrrrr}
 $\sigma_k/m$ & $L$ & $B_C$ & $q_{8,\min}$ & $\gamma_2$ \\
\colrule
0.05 & 24 & +0.528073 & 0.529525 & +0.477418 \\
0.08 & 12 & +0.437478 & 0.444632 & +0.367941 \\
0.08 & 16 & +0.293838 & 0.328346 & -1.699812 \\
0.08 & 24 & -0.071651 & 0.201665 & -4.584437 \\
0.10 & 6 & +0.539342 & 0.540250 & +0.521080 \\
0.10 & 8 & +0.425626 & 0.433961 & +0.409686 \\
0.10 & 12 & +0.194773 & 0.270423 & -2.398336 \\
0.10 & 16 & -0.109403 & 0.168785 & -4.746728 \\
0.15 & 3 & +0.512922 & 0.514393 & +0.498379 \\
0.15 & 4 & +0.391381 & 0.403673 & +0.372641 \\
0.15 & 6 & +0.131619 & 0.233573 & -0.060163 \\
0.15 & 8 & -0.225228 & 0.031703 & -0.335393 \\
0.20 & 2 & +0.467238 & 0.470767 & +0.411731 \\
0.20 & 3 & +0.260776 & 0.298145 & +0.240227 \\
0.20 & 4 & +0.011023 & 0.178925 & -0.047168 \\
0.30 & 1 & +0.471370 & 0.473952 & +0.464464 \\
0.30 & 2 & +0.008404 & 0.161144 & -0.767694 \\
0.40 & 0.5 & +0.584146 & 0.584312 & +0.578576 \\
0.40 & 1 & +0.224277 & 0.269906 & +0.196749 \\
0.50 & 0.5 & +0.441320 & 0.445665 & +0.374005 \\
0.50 & 1 & -0.111793 & 0.105516 & -0.138762 \\
\end{tabular}
\end{ruledtabular}
\end{table}



% ============================================================
\section{Static nuisance-domain sensitivity}
\label{sec:supp-K-sensitivity}
% ============================================================

On the Haberzettl slice $\delta=0$, the static protected fingerprint is
\begin{equation}
(W_M,W_C)=(g_M^2,1-g_M).
\end{equation}
For a compact nuisance interval $K$, define the uniform single-channel and combined Euclidean margins
\begin{equation}
\Delta_M(K)=\inf_{g\in K}g^2,
\qquad
\Delta_C(K)=\inf_{g\in K}|1-g|,
\end{equation}
\begin{equation}
\Delta_{MC}(K)=\inf_{g\in K}\sqrt{g^4+(1-g)^2}.
\end{equation}
These quantities give the exact static sensitivity to the admitted $g_M$ domain on the $\delta=0$ slice. The magnetic channel loses uniform separation exactly when $0\in K$, while the charge channel loses it exactly when $1\in K$. One protected primitive response therefore suffices whenever either blind point is excluded. Two are required when both are admitted. No numerical optimization is needed for this result.

Table~\ref{tab:K-sensitivity} evaluates representative intervals, including neighborhoods of the natural tree-level value $g_M=2$ and the broad interval used in the finite-bandwidth stress test. The combined margin remains positive throughout. Its global minimum occurs at $g=0.5897545123\ldots$ and equals $0.5378414487\ldots$, which is the closest approach in the chosen Euclidean coordinates.

\begin{table}[h]
\caption{Static nuisance-domain sensitivity. The last column is the minimum primitive protected support for uniform separation.}
\label{tab:K-sensitivity}
\begin{ruledtabular}
\begin{tabular}{lrrrr}
$K$ & $\Delta_M$ & $\Delta_C$ & $\Delta_{MC}$ & $\tau$\\
\colrule
$\{2\}$ & $4.000$ & $1.000$ & $4.123$ & 1\\
$[1.9,2.1]$ & $3.610$ & $0.900$ & $3.720$ & 1\\
$[1.5,2.5]$ & $2.250$ & $0.500$ & $2.305$ & 1\\
$[1,3]$ & $1.000$ & $0$ & $1.000$ & 1\\
$[0.5,3.5]$ & $0.250$ & $0$ & $0.538$ & 1\\
$[0,4]$ & $0$ & $0$ & $0.538$ & 2\\
$[-4,4]$ & $0$ & $0$ & $0.538$ & 2\\
\end{tabular}
\end{ruledtabular}
\end{table}

For the centered family $K_R=[2-R,2+R]$, the charge blind point enters at $R=1$ and the magnetic blind point enters at $R=2$. Therefore
\begin{equation}
\tau(K_R)=1\quad (0\le R<2),
\qquad
\tau(K_R)=2\quad (R\ge2).
\end{equation}
The corresponding static margins are plotted in the main article. The static calculations can also be reproduced from the project repository.

\subsection{Finite-width magnetic separation for $K=[1,3]$}
For the compact interval $K=[1,3]$, the magnetic anchor obeys $|g_M|\ge1$. The positive normalized magnetic measure and the Lipschitz condition give
\begin{equation}
\overline W_M\ge \int\dd\nu_M(x)\,[\max(1-Lx,0)]^2
\equiv M_L(1)>0.
\label{eq:supp-narrow-magnetic}
\end{equation}
The inequality holds for the whole Lipschitz class, independently of the charge form factor. Strict positivity follows because the integrand is positive for $0\le x<1/L$ and every finite-width Gaussian magnetic measure assigns positive mass to that interval. At $L=0$ the bound is one by continuity. The result is therefore a one-response certificate for every $0<\sigma_k/m<\infty$ and finite $L\ge0$ under this prior. The numerical grid reports the size of the analytic margin. The broad-class origin overlaps require $g_M=0$ and are excluded by $K=[1,3]$.

Table~\ref{tab:supp-narrow-magnetic} gives representative values from the $9\times11$ evaluation. All 99 margins are positive, with the smallest $0.005136151$ at $(\sigma_k/m,L)=(0.50,24)$. They are obtained by direct quadrature of the positive magnetic measure. The code and complete results are available in \path{theory_numerics/results/narrow_prior_magnetic_margins.csv} and \path{theory_numerics/narrow_prior_magnetic_audit.py}.
\begin{table}[h]
\caption{Finite-width magnetic-only lower bound $M_L(1)$ for $K=[1,3]$. This bound applies uniformly over all admitted form factors for each displayed $(\sigma_k/m,L)$.}
\label{tab:supp-narrow-magnetic}
\begin{ruledtabular}
\begin{tabular}{crr}
$\sigma_k/m$ & $L=2$ & $L=24$\\
\colrule
$0.02$ & $0.995214$ & $0.943831$\\
$0.10$ & $0.888252$ & $0.283866$\\
$0.20$ & $0.635442$ & $0.057708$\\
$0.50$ & $0.160533$ & $0.005136$\\
\end{tabular}
\end{ruledtabular}
\end{table}

For the centered family $K_R=[2-R,2+R]$ with $0\le R<2$, the same argument gives $\overline W_M\ge M_L(2-R)>0$. At $R=2$ the magnetic blind point becomes admissible, so the broad-class joint-response analysis is needed. We use $[1,3]$ as a theory-motivated neighborhood of the tree-level point, without treating it as a universal empirical prior for a new particle.

\subsection{Robustness to $g_M+g_Q=1+\delta$}
For a symmetric deformation domain $|\delta|\le\Delta$, magnetic positivity still reduces a possible protected-origin overlap to $g_M=0$ and $G_M\equiv0$. The remaining anchors satisfy $G_1(0)=1$ and $G_Q(0)=1+\delta$. A uniform pointwise relaxation therefore uses
\begin{equation}
h_{L,\Delta}(x)=
\min_{\substack{a\in[1-Lx,1+Lx]\\
b\in[1-\Delta-Lx,1+\Delta+Lx]}}
\left(ab+\frac{x}{4}b^2\right),
\end{equation}
with $B_C(\sigma_k/m,L,\Delta)=\int\dd\nu_C\,h_{L,\Delta}$. A positive value excludes a joint origin overlap uniformly over the complete enlarged nuisance class.

Table~\ref{tab:delta-robustness} shows the resulting sensitivity. The $33$ explicit overlap witnesses remain admissible at every displayed $\Delta$ because each enlarged class contains the $\delta=0$ slice. Points that lose the sufficient whole-class bound and have no explicit witness are reported as unresolved. The compact-prior result in Sec.~2.1 is unchanged because it uses only $G_M$.

\begin{table}[h]
\caption{Uniform robustness to the deformation $|\delta|\le\Delta$. The counts use the same 99-point bandwidth--smoothness grid. $B_C^{\rm ref}$ and $D^{\rm ref}$ are the charge-envelope and unit-covariance distance lower bounds at $(\sigma_k/m,L)=(0.10,2)$.}
\label{tab:delta-robustness}
\begin{ruledtabular}
\begin{tabular}{rrrrrr}
$\Delta$ & certified & overlap & unresolved & $B_C^{\rm ref}$ & $D^{\rm ref}$\\
\colrule
0.00 & 63 & 33 & 3 & +0.828236 & 0.403483\\
0.05 & 61 & 33 & 5 & +0.782066 & 0.369644\\
0.10 & 61 & 33 & 5 & +0.735957 & 0.336361\\
0.25 & 59 & 33 & 7 & +0.597992 & 0.240405\\
0.50 & 55 & 33 & 11 & +0.369261 & 0.099080\\
0.75 & 45 & 33 & 21 & +0.141370 & 0.009804\\
0.90 & 31 & 33 & 35 & -0.002478 & 0.000000\\
\end{tabular}
\end{ruledtabular}
\end{table}

At the reference point, the uniform charge envelope remains positive until $\Delta\simeq0.897582$. This number characterizes the stated pointwise whole-class relaxation and is not an empirical bound on an anomalous coupling. In the standard CP-even $WW\gamma$ convention, $\delta=2\lambda_\gamma$ \cite{HagiwaraPecceiZeppenfeldHikasa1987}. The static protected distance remains positive for every $\Delta<1$ and reaches zero when the allowed domain first contains the common blind point $(g_M,\delta)=(0,-1)$.

% ============================================================
\section{Covariance-shape sensitivity}
% ============================================================

The physical finite-noise metric is
\begin{equation}
\Delta_\Sigma=\sqrt{\Delta\bm W^{\mathsf T}\Sigma^{-1}\Delta\bm W},
\qquad
\Sigma=\varepsilon^2
\begin{pmatrix}
1 & \rho r_\Sigma\\
\rho r_\Sigma & r_\Sigma^2
\end{pmatrix},
\end{equation}
where $r_\Sigma=\sigma(W_C)/\sigma(W_M)$. Table~\ref{tab:covariance} gives the whole-class lower bound at the reference point $\sigma_k/m=0.10$, $L=2$ for the tested covariance shapes. The numerical finite-noise gap depends on the detector covariance shape. With this parametrization $\varepsilon=\sigma(W_M)$ and $\sigma(W_C)=r_\Sigma\varepsilon$. The scale $\varepsilon$ becomes a common standard deviation only when $r_\Sigma=1$.

\begin{table}[h]
\caption{Covariance-shape sensitivity at $\sigma_k/m=0.10$ and $L=2$. The lower bound follows from the whole-class half-strip enclosure, not only from exact-overlap exclusion. The last two columns give the magnetic-coordinate reference scale allowed for nominal $3\sigma$ and $5\sigma$ separation.}
\label{tab:covariance}
\begin{ruledtabular}
\begin{tabular}{ccccc}
$r_\Sigma$ & $\rho$ & lower bound & $\varepsilon_{3\sigma}$ & $\varepsilon_{5\sigma}$\\
\colrule
0.5 & 0 & 0.5541 & 0.1847 & 0.1108\\
0.5 & 0.5 & 0.4564 & 0.1521 & 0.0913\\
0.5 & -0.5 & 0.7689 & 0.2563 & 0.1538\\
1 & 0 & 0.4035 & 0.1345 & 0.0807\\
1 & 0.5 & 0.3299 & 0.1100 & 0.0660\\
1 & -0.5 & 0.5691 & 0.1897 & 0.1138\\
2 & 0 & 0.2579 & 0.0860 & 0.0516\\
2 & 0.5 & 0.2150 & 0.0717 & 0.0430\\
2 & -0.5 & 0.3532 & 0.1177 & 0.0706\\
4 & 0 & 0.1506 & 0.0502 & 0.0301\\
4 & 0.5 & 0.1293 & 0.0431 & 0.0259\\
4 & -0.5 & 0.1986 & 0.0662 & 0.0397\\
\end{tabular}
\end{ruledtabular}
\end{table}

% ============================================================
\section{Estimator-space recovery}
% ============================================================

Statistical recovery is evaluated directly in calibrated fingerprint space using estimator uncertainty. The estimator is sampled as
\begin{equation}
\widehat{\bm W}\sim\mathcal N(\bm W,\varepsilon^2\Sigma_{\rm shape}).
\end{equation}
At $\sigma_k/m=0.10$, $L=2$, and $\Sigma_{\rm shape}=I$, the profile-distance statistic is calibrated to a $5\%$ false-positive rate under spin $1/2$. The 30000-point inference cloud draws $g\sim\mathrm{Unif}[-4,4]$ and each of the 24 slopes of $(G_1,G_M,G_Q)$ independently from $\mathrm{Unif}[-2,2]$. The 10000-point truth cloud is drawn independently from the same distribution. The first power column in Table~\ref{tab:power} is therefore nuisance averaged. The second power column uses the explicit admissible profile with $g=0.5734946$ and $(\overline W_M,\overline W_C)=(0.2672378,0.3023109)$, which supplies the upper endpoint of the tight class-distance bracket.

\begin{table}[h]
\caption{Estimator-space spin-$1$ recovery at the reference finite-bandwidth point. ``Cloud'' denotes the stated nuisance average. ``Near boundary'' denotes the fixed explicit bracketing profile. The uncertainty scale $\varepsilon$ is an estimator-space scale.}
\label{tab:power}
\begin{ruledtabular}
\begin{tabular}{cccc}
$\varepsilon$ & power, cloud & power, near boundary & AUC, cloud\\
\colrule
0.50 & 0.881 & 0.166 & 0.9690\\
0.30 & 0.961 & 0.357 & 0.9916\\
0.20 & 0.992 & 0.629 & 0.9982\\
0.15 & 0.999 & 0.824 & 0.9997\\
\end{tabular}
\end{ruledtabular}
\end{table}


% ============================================================
\section{Ablation, finite-basis checks, and the Standard Model point}
% ============================================================

\subsection{Primitive-resource ablation}
At $\sigma_k/m=0.10$ and $L=2$, choosing $g_M=G_M\equiv0$ gives an exact magnetic-only blind point. The charge-only response vanishes for $G_Q\equiv0$ at $g_M=1$ on $\delta=0$. The independently derived two-channel unit-covariance bracket is $0.403483\le d_{\rm class}\le0.403495$. The finite-width ablation preserves the broad-class static structure. Each primitive coordinate can be tuned away separately, while the pair remains separated at the reference point.

\subsection{Expanded moment-relation negative control}
In the deformation $g_M+g_Q=1+\delta$, the point $g_M=0$ and $\delta=-1$ gives $g_Q=0$ and therefore reaches the protected origin. This is the first common blind point of the two primitive responses. The control isolates the role of the moment-relation nuisance class from the representation-theoretic fact that spin $1$ permits a rank-$2$ sector.

\subsection{Finite-basis and whole-class comparison}

At $(\sigma_k/m,L)=(0.15,8)$, the global eight-segment QCQP minimum is positive even though an admissible two-kink overlap exists. A finite basis can therefore miss a degeneracy. The $48$-segment multistart search recovers a negative profile at this point. For $(0.20,2)$, the analytic pointwise bound gives $B_C=0.467238>0$ and certifies the complete anchored Lipschitz class independently of the basis size. The numerical searches are used to locate admissible counterexamples and to check the optimization procedure.

\subsection{Standard Model $W$ point}
On the reference $\delta=0$ slice, the scan includes $g_M=2$, so the static relation gives $g_Q=-1$. The corresponding protected static fingerprint is
\begin{equation}
\bm W_{\rm YM}=(4,-1),
\end{equation}
with unit-covariance distance $\sqrt{17}$ from the spin-$1/2$ origin. This point is a control that the nuisance range contains the canonical charged Yang--Mills vector normalization \cite{HagiwaraPecceiZeppenfeldHikasa1987,Haberzettl2019,BrodskyHiller1992}.

% ============================================================
\section{Protected-sector extension to other spins}

Let $s<s'$ and let $\Pi_{>2s}$ project onto the internal rotational tensor ranks unavailable to spin $s$. For a specified higher-spin nuisance set $\Xi_{s'}$, the protected response set is
\begin{equation}
\mathcal V_{s'|s}=\{\Pi_{>2s}\bm W_{s'}(\xi):\xi\in\Xi_{s'}\}\subseteq\mathcal R_{\mathcal A}.
\end{equation}
The lower-spin model has zero response in this protected sector. For a primitive detector block $a\in\mathcal A$, define its blind locus
\begin{equation}
Z_a=\{\xi\in\Xi_{s'}:\pi_a\Pi_{>2s}\bm W_{s'}(\xi)=0\}.
\end{equation}
The architecture-aware separation threshold is therefore
\begin{equation}
\tau_{\mathcal A}(s,s')=\min_{S\subseteq\mathcal A}\left\{|S|:\bigcap_{a\in S}Z_a=\varnothing\right\}.
\end{equation}
It counts the primitive response blocks needed to exclude a common blind point in the admitted higher-spin nuisance class. Its use requires a separately justified electromagnetic multipole model for that class. For finite polynomial parametrizations, moment and sum-of-squares methods can supply additional sufficient bounds \cite{Lasserre2001,Parrilo2003}. Positive finite-basis certificates need not establish separation for an unrestricted functional class.


\section{Reproducibility}
% ============================================================
The $99$-point classification, overlap witnesses, distance bounds, covariance analysis, QCQP/SOS results, estimator-space calculations and $\delta$-robustness tests are available at the fixed repository commit cited in the main article. The $48$-segment search code, saved slope profiles, independent quadrature checks and convergence table are available in the same fixed repository snapshot. The primary data are in \texttt{theory\_numerics/results/alpha0/}.

\FloatBarrier
\section*{References}
\printbibliography[heading=none]
\end{document}
